\section{Moment matrices and higher-rank amplituhedral realizations}
\label{sec:higher-rank}
The partition of the dodecaphasic record extends the rank-one
construction to explicit higher-rank points. The same thirteen
ordered nodes simultaneously provide positive Plücker minors
and a positive-definite moment form. These two positivities are
obtained from a single list of intervals; the second proves that the
composition preserves the required rank.

Let \(d_j=K_j/Q\), where \(Q=\sum_{j=1}^{12}K_j=6263\), and define
\[
 t_0=0,\qquad t_j=\sum_{r=1}^jd_r.
\]
Positivity of the twelve blocks gives
\(0=t_0<t_1<\cdots<t_{12}=1\). For each \(2\le k\le9\), set
\[
 C^{(k)}_{ai}=t_i^a\quad(0\le a<k),\qquad
 Z^{(k)}_{bi}=t_i^b\quad(0\le b<k+4),\qquad 0\le i\le12.
\]
The convention \(t_i^0=1\) includes the zero node. Thus, \(C^{(k)}\)
is a \(k\times13\) matrix, while \(Z^{(k)}\) has
size \((k+4)\times13\).

\begin{theorem}[Positivity and rank of the moment realization]
The matrix \(C^{(k)}\) represents a point of
\(\Gr_{>0}(k,13)\), and all ordered maximal minors of
\(Z^{(k)}\) are positive. The matrix
\[
 Y^{(k)}=C^{(k)}(Z^{(k)})^{\mathsf T},\qquad
 Y^{(k)}_{ab}=\sum_{i=0}^{12}t_i^{a+b},
\]
has rank \(k\). In particular, for the positive external datum
\(Z^{(k)}\) already constructed from the record,
\[
 [Y^{(k)}]\in\mathcal A_{13,k,4}(Z^{(k)}),\qquad 2\le k\le9,
\]
where
\[
 \mathcal A_{13,k,4}(Z^{(k)})
 =\{[C(Z^{(k)})^{\mathsf T}]:[C]\in\Gr_{\ge0}(k,13)\}.
\]
\end{theorem}
\begin{proof}
A maximal minor of either initial matrix has the form
\[
 \det(t_{i_j}^{a})_{a=0}^{r-1}
 =\prod_{p<q}(t_{i_q}-t_{i_p})>0,
\]
for increasing indices. The inequality proves both Plücker
assertions. The initial \(k\times k\) block of \(Y^{(k)}\) is the
Hankel matrix \(H^{(k)}\). For \(x\ne0\),
\[
 x^{\mathsf T}H^{(k)}x
 =\sum_{i=0}^{12}\left(\sum_{a=0}^{k-1}x_at_i^a\right)^2>0.
\]
Indeed, a nonzero polynomial of degree less than \(k\le9\) cannot
vanish at thirteen distinct points. Hence \(H^{(k)}\) is positive definite,
its determinant is strictly positive and \(Y^{(k)}\) has rank
\(k\). Amplituhedral membership is then the composition of the
two preceding positive matrices, with rank already verified.
\end{proof}

In rank two, the proof takes the particularly explicit form
\[
 \det H^{(2)}=13\sum_i t_i^2-\left(\sum_i t_i\right)^2
 =\sum_{i<j}(t_j-t_i)^2>0.
\]
Node separation, produced by positivity of the record,
is thus converted into coercivity of a quadratic form. Each
fixed record corresponds to a specified point for each rank. The
complete coverage, triangulations and residues proved in the
preceding section retain their rank-one domain; the present theorem
adds the explicit collection of rank-two to rank-nine points and their moment
matrices. This distinction allows each proof to be reused within its
precise scope.
